\documentclass[11pt]{article}
\usepackage[margin=1in]{geometry}
\usepackage{amsmath,amssymb,amsthm}
\usepackage{mathtools}
\usepackage{hyperref}
\usepackage{xcolor}
\usepackage{enumitem}

\newtheorem{theorem}{Theorem}
\newtheorem{proposition}{Proposition}
\newtheorem{lemma}{Lemma}

\newcommand{\bbY}{Y}
\newcommand{\bbS}{S}
\newcommand{\bbU}{U}
\newcommand{\bbX}{X}
\newcommand{\bbZ}{Z}
\newcommand{\bbV}{V}
\newcommand{\E}{\mathbb{E}}
\newcommand{\BigOh}{\mathcal{O}}

\title{From Single-User SC Trellis Decoding to a Two-User MAC \\
       Successive Cancellation Trellis (SCT) Decoder}
\author{Polar Codes for MAC Project --- Implementation Notes}
\date{\today}

\begin{document}
\maketitle

\begin{abstract}
We summarise the single-user successive cancellation (SC) trellis decoder of
Wang, Honda, Yamamoto, Liu and Hou (ITW 2015) for polar codes over
finite-state channels with memory, and then derive its two-user
multiple-access channel (MAC) extension by composing it with the monotone
chain rule decomposition of \"{O}nay (ISIT 2013).  The two-user version is
the \emph{SCT decoder} used as the analytical baseline in our project.  We
state the recursive formulas, give a one-paragraph justification of why they
are correct, and describe how the algorithm is realised in our repository.
\end{abstract}

%-----------------------------------------------------------------------------
\section{Setting}

A discrete finite-state channel is specified by
\[
  W\bigl(y_i, s_i \mid x_i, s_{i-1}\bigr),
\]
where $x_i \in \mathcal{X}=\{0,1\}$ is the input, $y_i \in \mathcal{Y}$ the
output, and $s_i \in \mathcal{S}$ a finite hidden state evolving in time.
The state at time $0$ is known (say $s_0$).  Memoryless channels are the
special case $|\mathcal{S}|=1$.

\paragraph{Polar transform.}
Let $\mathbf{G}_N = G^{\otimes n}$ with $G=\begin{psmallmatrix}1&0\\1&1\end{psmallmatrix}$,
$N=2^n$.  The input is $\bbX_1^N = \bbU_1^N \mathbf{G}_N$; SC decoding aims to recover
$\bbU_1^N$ from $\bbY_1^N$.  S\c{a}so\u{g}lu (2011) showed that the polar
transform polarises finite-state channels in exactly the same way it
polarises B-DMCs.  What is needed in practice is a \emph{decoder} that runs
in time polynomial in $|\mathcal{S}|$ and $N$.

%-----------------------------------------------------------------------------
\section{Single-User SCT (Wang/Honda/Yamamoto/Liu/Hou, ITW 2015)}

\subsection{Augmented joint probability}

For memoryless channels the SC decoder uses
\[
  W_N^{(i)}\bigl(y_1^N, u_1^i\bigr) := P\bigl(\bbY_1^N=y_1^N, \bbU_1^i=u_1^i\bigr),
\]
which factorises along the polar tree.  With memory this no longer
factorises because the channel has correlated past and future.  The paper's
key idea is to add the \emph{boundary states} $s_0, s_N$ to the conditioning:
\begin{equation}\label{eq:augmented}
  W_N^{(i)}\bigl(y_1^N, u_1^i, s_N \,\big|\, s_0\bigr)
  := P\bigl(\bbY_1^N=y_1^N,\; \bbU_1^i=u_1^i,\; \bbS_N=s_N \,\big|\, \bbS_0=s_0\bigr).
\end{equation}
Conditioning on the entry state $s_0$ and pinning the exit state $s_N$ makes the
augmented quantity factor cleanly over a half-block, because everything
before $s_N$ is conditionally independent of everything after it given $s_N$.

\subsection{Recursive formulas}

The paper proves (Theorem 2 in \cite{Wang2015}) that for $N=2^n$,
$1 \le i \le N$,
\begin{align}
  W_{2N}^{(2i-1)}\bigl(y_1^{2N}, u_1^{2i-1}, s_{2N} \mid s_0\bigr)
  &= \sum_{s_N}\sum_{u_{2i}}
     W_N^{(i)}\bigl(y_1^N,\; u_{1,o}^{2i}\oplus u_{1,e}^{2i},\; s_N \mid s_0\bigr)
     \;\cdot\; \nonumber\\
  &\hspace{2.4cm}
     W_N^{(i)}\bigl(y_{N+1}^{2N},\; u_{1,e}^{2i},\; s_{2N}\mid s_N\bigr),
     \label{eq:su-odd}\\[4pt]
  W_{2N}^{(2i)}\bigl(y_1^{2N}, u_1^{2i}, s_{2N} \mid s_0\bigr)
  &= \sum_{s_N}
     W_N^{(i)}\bigl(y_1^N,\; u_{1,o}^{2i}\oplus u_{1,e}^{2i},\; s_N \mid s_0\bigr)
     \;\cdot\; \nonumber\\
  &\hspace{2.4cm}
     W_N^{(i)}\bigl(y_{N+1}^{2N},\; u_{1,e}^{2i},\; s_{2N}\mid s_N\bigr).
     \label{eq:su-even}
\end{align}
Here $u_{1,o}^{2i}$ and $u_{1,e}^{2i}$ are the odd- and even-indexed
sub-sequences of $u_1^{2i}$ --- the standard polar split.

\subsection{Why the recursion is correct}

The argument has three ingredients:
\begin{enumerate}[leftmargin=*]
\item The polar transform splits $u_1^{2i}$ into
      $\bigl(u_{1,o}\oplus u_{1,e}\bigr) \mathbf{G}_N$ in the first half-block and
      $u_{1,e}\mathbf{G}_N$ in the second --- exactly as for memoryless channels.
\item Conditional on $s_N$, the past $\bigl(\bbY_1^N, \bbX_1^N, s_0\bigr)$
      is independent of the future $\bigl(\bbY_{N+1}^{2N}, \bbX_{N+1}^{2N}, s_{2N}\bigr)$.
      So the joint factorises, and the two half-blocks each contribute a
      single $W_N^{(i)}$ term.
\item The difference between the odd index ($2i-1$) and the even index
      ($2i$) is exactly Arikan's split: the odd case still has to marginalise
      $u_{2i}$, the even case has already absorbed it.
\end{enumerate}
Bringing those together gives \eqref{eq:su-odd}--\eqref{eq:su-even}.  All
inner steps reduce to ordinary marginalisation.

\subsection{Complexity}

The recursion is the same polar tree as the memoryless case, but at each of
the $\BigOh(N\log N)$ steps it must sum over $|\mathcal{S}|$ choices of
$s_N$ and store a value for every $(s_0, s_N) \in \mathcal{S}^2$.  Total
cost $\BigOh\bigl(|\mathcal{S}|^3 N \log N\bigr)$.  For ISI with a 1-tap
state, $|\mathcal{S}|=2$, so this stays $\BigOh(N\log N)$ up to constants.

\subsection{Recovering the per-bit posterior}

The decoder needs the posterior $P\bigl(\bbU_i \mid \bbU_1^{i-1}, \bbY_1^N\bigr)$.
From $W_N^{(i)}$ the joint probability
$P(\bbY_1^N, \bbU_1^i)$ is obtained by averaging over the boundary states,
\begin{equation}\label{eq:marginal}
  P\bigl(\bbY_1^N=y_1^N,\; \bbU_1^i=u_1^i\bigr)
  = \sum_{s_0,\,s_N} W_N^{(i)}(y_1^N, u_1^i, s_N\mid s_0)\;P(\bbS_0=s_0),
\end{equation}
and the SC decision rule is the usual $\hat{u}_i = \arg\max_{u_i}
P(\bbY_1^N, \bbU_1^{i-1}\hat{u}_1^{i-1}, u_i)$.

%-----------------------------------------------------------------------------
\section{Lifting to the Two-User MAC}\label{sec:lift}

\subsection{Two-user MAC setting}

A two-user binary-input MAC with memory has transition law
\[
  W\bigl(z_i, s_i \mid x_i, y_i, s_{i-1}\bigr),
\]
where $x_i$ and $y_i$ are the two users' inputs, $z_i$ the joint output, and
$s_i$ the channel state.  The state $s_i$ now generally has to encode all
information needed to predict the next output.  For the ISI--MAC
\[
  Z_i = (1-2X_i) + (1-2Y_i) + h\bigl[(1-2X_{i-1}) + (1-2Y_{i-1})\bigr] + W_i,
  \qquad W_i\sim\mathcal{N}(0,\sigma^2),
\]
the natural state is $s_i = (X_{i-1}, Y_{i-1}) \in \{0,1\}^2$, so
$|\mathcal{S}| = 4$.  Inputs are independent uniform: $P_{X}(0)=P_{X}(1)=
P_{Y}(0)=P_{Y}(1)=\tfrac{1}{2}$.

\subsection{\"{O}nay's monotone-chain decomposition (ISIT 2013)}

Let $\bbU^N$ and $\bbV^N$ be the polar pre-codewords for the two users.
\"{O}nay~\cite{Onay2013} expands $I(\bbZ^N;\bbU^N,\bbV^N)$ along a
\emph{monotone chain rule path} $b^{2N}\in\{0,1\}^{2N}$: a $0$ at step $k$
emits an extra $\bbU$ symbol, a $1$ emits an extra $\bbV$ symbol.  Each
monotone path produces a sequence of \emph{coordinate channels}
$W_N^{(b_k,i,j)}$, $k=i+j$, that are single-user binary-input channels
(eq.~(8) of \cite{Onay2013}):
\begin{equation}\label{eq:onay-coord}
  W_N^{(b_k,i,j)}(z^N, s^{k-1} \mid s_k)
  =\begin{cases}
     W_N^{(0,i,j)}(z^N, u^{i-1}, v^j \mid u_i) & \text{if } b_k=0,\\
     W_N^{(1,i,j)}(z^N, u^i, v^{j-1} \mid v_j) & \text{if } b_k=1,
   \end{cases}
\end{equation}
and these polarise (capacity terms go to $0$ or $1$) for every monotone
path. The decoder works through $k=1,\ldots,2N$ deciding one bit at a step,
of user-1 if $b_k=0$ or user-2 if $b_k=1$.  In particular:
\begin{itemize}
\item $b^{2N}=0^N 1^N$ is the \emph{corner-rate Class C} path:
      decode all of $\bbU$ first treating $\bbV$ as i.i.d.\ noise, then all
      of $\bbV$ conditioned on the just-decoded $\hat{\bbX} = \hat{\bbU}\mathbf{G}_N$.
\item Other paths give intermediate rate splits and reach the dominant face
      of the MAC capacity region with arbitrary precision.
\end{itemize}
\"{O}nay's recursive formulas (his eqs.~(16)--(19)) are simply the same polar
recursion as in the single-user case, but they propagate the four-type
coordinate channels $W_N^{(b_k,i,j)}$ instead of one type.

\subsection{Composition with the state-lattice trick}

Both \"{O}nay's MAC SC and Wang/Honda et al.'s state-lattice SC use exactly
the same polar split.  The composition is therefore straightforward:
augment \"{O}nay's coordinate channels with the boundary states $(s_0, s_N)$
exactly as in \eqref{eq:augmented}.  Define
\begin{equation}\label{eq:mac-augmented}
  W_N^{(b_k,i,j)}\bigl(z_1^N, u^{i-1}\text{ or }u^i,\; v^{j-1}\text{ or }v^j,\;
                        s_N \,\big|\, u_i \text{ or } v_j,\; s_0\bigr),
\end{equation}
keeping all four \"{O}nay types $b_k=0$ vs $b_k=1$, even vs odd index.  The
recursion for each type takes exactly the same shape as the single-user
SCT recursion \eqref{eq:su-odd}--\eqref{eq:su-even}: a sum over the
shared boundary state $s_N$ inside the product of two half-block factors.
For example, the analogue of \eqref{eq:su-odd} for an odd-index
$b_k=0$ ($\bbU$) step is
\begin{align}
   &W_{2N}^{(0, 2i-1, 2j)}\bigl(z_1^{2N}, u_1^{2i-2}, v_1^{2j},\; s_{2N} \,\big|\, u_{2i-1},\; s_0\bigr) \nonumber\\
  &\hspace{1.2cm} = \frac{1}{2^2}\sum_{s_N}\;\sum_{u_{2i}, v_{2j-1}, v_{2j}}\;
        W_N^{(0,i,j)}\bigl(z_1^N,\, \tilde u^{i-1},\, \tilde v^{j},\, s_N \,\big|\, \bar u_i,\, s_0\bigr) \nonumber\\
  &\hspace{4.8cm}\cdot
        W_N^{(0,i,j)}\bigl(z_{N+1}^{2N},\, \bar u^{i-1},\, \bar v^{j},\, s_{2N} \,\big|\, \bar u_i,\, s_N\bigr),
   \label{eq:mac-recursion}
\end{align}
where $\tilde u = u_{1,o}\oplus u_{1,e}$, $\bar u = u_{1,e}$ and similarly
for $\tilde v$, $\bar v$.  Three other shape-equivalent identities cover the
$(b_k, \mathrm{odd/even})\in\{0,1\}\times\{\text{odd,even}\}$ cases; they
are the natural composition of \"{O}nay's eqs.~(16)--(19) with the
state-summation of \eqref{eq:su-odd}--\eqref{eq:su-even}.  No new
information-theoretic argument is needed beyond:
\begin{enumerate}[leftmargin=*]
\item the polar split applies to $u$ and $v$ separately (\"{O}nay), and
\item conditional on $s_N$, the past and future of the channel are
      independent (Wang/Honda).
\end{enumerate}

\subsection{Complexity}

Each step of the recursion for the MAC version sums over $|\mathcal{S}|$
choices of $s_N$ and stores a table indexed by $(s_0, s_N)$, exactly as
single-user.  Thus the cost remains
\(\BigOh(|\mathcal{S}|^3 N \log N)\).
For ISI--MAC, $|\mathcal{S}|=4$ and the algorithm runs in
$\BigOh(64\, N\log N) = \BigOh(N\log N)$.

%-----------------------------------------------------------------------------
\section{Corner-Rate Specialisation and Equivalent Two-Stage Form}

For the corner-rate path $b^{2N}=0^N 1^N$, all $\bbU$ steps come before any
$\bbV$ step.  The recursion \eqref{eq:mac-recursion} simplifies because, in
the $\bbU$ phase, each MAC coordinate channel $W^{(0,i,j)}$ with $j=N$
\emph{marginalises} $\bbV$ as uniform i.i.d.: the joint channel
$W(z\mid x,y)$ is replaced by its $\bbV$-averaged version
\(
   \tilde W(z\mid x) = \tfrac{1}{2}\sum_y W(z\mid x, y).
\)
Equivalently, the entire $\bbU$ phase can be implemented in two steps:
\begin{enumerate}[leftmargin=*]
\item[\textbf{(a)}] Compute the per-position marginal
      $\log P(z_t \mid x_t)$ from \emph{joint} forward--backward on the
      $|\mathcal{S}|$-state trellis with i.i.d.\ uniform $\bbY$.  This
      produces a scalar LLR $\ell^{(u)}_t$ for each position.
\item[\textbf{(b)}] Run \emph{ordinary} Arikan SC on $\ell^{(u)}_1, \ldots,
      \ell^{(u)}_N$ to decide $\hat{\bbU}$.
\end{enumerate}
For stage 2, condition on $\hat{\bbX} = \hat{\bbU}\mathbf{G}_N$: now both $X_t$ and
$X_{t-1}$ are known, so the state collapses to $s_t = Y_{t-1}$, a 2-state
chain.  Run a 2-state FB to obtain a scalar LLR $\ell^{(v)}_t$ and
ordinary Arikan SC for $\hat{\bbV}$.  This is the version used in the
project; it is mathematically equivalent to \eqref{eq:mac-recursion}
specialised to $b^{2N}=0^N 1^N$.

%-----------------------------------------------------------------------------
\section{Realisation in the Project}

Two implementations follow the development above; both are validated to
agree within Monte-Carlo noise across $N\in\{16,\ldots,1024\}$:

\begin{itemize}[leftmargin=*]
\item \texttt{polar/decoder\_trellis.py} (\emph{joint MAC SCT}): builds the
  4-state leaf tensor for ISI--MAC, runs a single \texttt{forward-backward}
  on the full $|\mathcal{S}|=4$ lattice, and feeds the resulting
  $(N,2,2)$ per-position marginals into a Ren-style computational-graph
  two-user MAC SC tree.  Supports any monotone path $b^{2N}$ via
  \texttt{polar/decoder\_interleaved.py}.

\item \texttt{scripts/local\_analysis/chained\_sct\_4state.py} (\emph{4-state
  chained corner-rate}): the two-stage form of Section 4 above.  Stage 1
  uses the same 4-state FB but marginalises $\bbY$ per position to obtain
  scalar U-LLRs, then runs Arikan SC.  Stage 2 conditions on $\hat{\bbX}$
  and runs a 2-state $Y_{t-1}$ FB followed by Arikan SC.  This decoder is
  the primary analytical baseline we plot against the neural decoders.
\end{itemize}

\paragraph{Numerical sanity.}  At ISI--MAC with $h=0.3$, SNR $=6$ dB,
corner-rate path, the 4-state chained SCT matches NPD within
Monte-Carlo noise at every $N\in\{16,\ldots,1024\}$ once the MC
genie-aided design uses $\ge 50{,}000$ trials per $N$ (and
$\ge 200{,}000$ at $N=1024$ to resolve the long reliability tail).

%-----------------------------------------------------------------------------
\section{Summary}

\begin{itemize}[leftmargin=*]
\item The single-user SCT of Wang/Honda/Yamamoto/Liu/Hou (2015) generalises
      Arikan SC to finite-state channels by adding the boundary states
      $(s_0, s_N)$ to every polar-tree intermediate quantity, see
      \eqref{eq:augmented}.  The polar-tree recursion is unchanged in form
      \eqref{eq:su-odd}--\eqref{eq:su-even}; only an extra sum over $s_N$ is
      added at each step.
\item \"{O}nay (2013) showed that any monotone chain rule expansion of
      $I(Z; U, V)$ gives a polar code reaching every point on the dominant
      face of the MAC capacity region.  His decoder is the same polar-tree
      recursion as single-user SC, applied to the four coordinate channels.
\item Composing the two is mechanical: take \"{O}nay's recursions for the
      four coordinate-channel types, augment each by $(s_0, s_N)$ and add a
      sum over $s_N$ at every composition --- exactly as in Wang/Honda et
      al.  The result is the two-user MAC SCT used in the project,
      complexity $\BigOh(|\mathcal{S}|^3 N \log N)$.
\item For the corner-rate path $b^{2N}=0^N 1^N$, the algorithm specialises
      to the two-stage form ``state-lattice FB $\to$ scalar LLR $\to$
      Arikan SC'', which is what we implement and use as our analytical
      baseline against the NPD and NCG decoders.
\end{itemize}

\begin{thebibliography}{9}
\bibitem{Wang2015}
R.~Wang, J.~Honda, H.~Yamamoto, R.~Liu and Y.~Hou,
``Construction of Polar Codes for Channels with Memory,''
\emph{Proc.\ IEEE Information Theory Workshop (ITW), Fall 2015}, pp.~187--191.
\bibitem{Onay2013}
S.~\"{O}nay,
``Successive Cancellation Decoding of Polar Codes for the Two-User Binary-Input MAC,''
\emph{Proc.\ IEEE Int.\ Symp.\ Inform.\ Theory (ISIT)}, 2013, pp.~1122--1126.
\bibitem{Sasoglu2011}
E.~\c{S}a\c{s}o\u{g}lu,
``Polarization in the Presence of Memory,''
\emph{Proc.\ IEEE ISIT}, 2011, pp.~189--193.
\end{thebibliography}

\end{document}
